\label{sec:framework}

\subsection{Economic Character as an M-Track Community Signal}

The M-track framework (M.0) defines community character as a function
$\mathbf{c}: \mathcal{T} \to \mathbb{R}^k_{\geq 0}$ mapping each census
tract to a non-negative character vector \citep{m0framework}.
The community similarity between adjacent tracts $u$ and $v$ is
$\mathrm{sim}(\mathbf{c}_u, \mathbf{c}_v) \in [0, 1]$, and the edge weight
modifier is $w(u,v) = w_{\mathrm{base}}(u,v) \times \mathrm{sim}(\mathbf{c}_u, \mathbf{c}_v)$.

For M.1, the character function is the economic character vector:
\begin{equation}
  \mathbf{e}(t) = \bigl(\mathrm{CI}(t),\ \mathrm{IF}(t),\ \mathrm{JPR}(t)\bigr)
  \label{eq:econchar}
\end{equation}
where CI, IF, and JPR are defined in Section~\ref{sec:background}.

This is a plug-in implementation of the M.0 framework: the cosine similarity
of economic character vectors replaces the generic $\mathrm{sim}$ in the
M.0 edge weight formula.

\subsection{Prototypical Tract Types and Their Similarity}

To build intuition for the cosine similarity of economic character vectors,
Table~\ref{tab:prototypes} shows three prototypical tract types and
their character vectors.

\begin{table}[ht]
\centering
\caption{Prototypical economic character vectors for three tract types.
  CI = commercial intensity, IF = industrial fraction, JPR = jobs per resident.}
\label{tab:prototypes}
\begin{tabular}{llrrr}
\toprule
Type & Example & CI & IF & JPR \\
\midrule
Bedroom suburb & Cary, NC residential & 0.02 & 0.01 & 0.3 \\
Commercial corridor & Charlotte retail strip & 0.45 & 0.05 & 2.1 \\
Industrial park & RDU airport logistics & 0.08 & 0.62 & 4.8 \\
\bottomrule
\end{tabular}
\end{table}

The cosine similarities between these prototypes reveal the geometry of
economic community membership.
Table~\ref{tab:similarity-matrix} shows the pairwise similarity matrix.

\begin{table}[ht]
\centering
\caption{Cosine similarity between prototypical tract types.
  Values above 0.8 indicate strong community cohesion (METIS avoids cutting);
  values below 0.3 indicate natural community boundaries (METIS drawn to cut here).}
\label{tab:similarity-matrix}
\begin{tabular}{lrrr}
\toprule
 & Bedroom & Commercial & Industrial \\
\midrule
Bedroom suburb     & 1.00 & 0.11 & 0.04 \\
Commercial corridor & 0.11 & 1.00 & 0.28 \\
Industrial park    & 0.04 & 0.28 & 1.00 \\
\bottomrule
\end{tabular}
\end{table}

The similarity matrix encodes the economic community structure of the
modern metropolitan area:
\begin{itemize}
  \item \textbf{Bedroom to bedroom (0.98 for identical vectors).}
    Two residential suburbs are economically homogeneous.
    They belong in the same district from an economic community perspective.
    METIS strongly avoids cutting this edge.
  \item \textbf{Bedroom to commercial (0.11).}
    The boundary between a residential neighborhood and the commercial
    corridor serving it is a natural cut seam.
    The residents and the businesses are different communities: the residents
    are customers and neighbors; the businesses are employers with different
    policy interests (commercial property tax rates, zoning variance).
  \item \textbf{Bedroom to industrial (0.04).}
    An even harder natural cut.
    The residents and the industrial park share almost no economic character.
    District cuts at the residential-industrial transition are among the
    strongest economic community signals the algorithm can exploit.
  \item \textbf{Commercial to commercial (0.95 for similar vectors).}
    Two commercial districts belong together.
    They share customer bases, commercial zoning interests, and economic
    policy concerns (retail leakage, big-box competition, downtown
    revitalization).
  \item \textbf{Commercial to industrial (0.28).}
    A softer cut---commercial and industrial zones are both non-residential
    but have genuinely different economic characters.
    This boundary is a soft invitation to cut, not a hard cut seam.
  \item \textbf{Industrial to industrial (0.95+ for similar vectors).}
    Industrial communities belong together: they share zoning interests
    (industrial land preservation), logistics infrastructure needs, and
    employer-community relationships.
\end{itemize}

\subsection{The Blend Formula}

Following B.27, the edge weight is:
\begin{equation}
  w_{\mathrm{ec}}(u, v)
  = w_{\mathrm{geo}}(u, v)
  \times \bigl(\alpha + (1 - \alpha)\,\mathrm{sim}(\mathbf{e}_u, \mathbf{e}_v)\bigr)
  \label{eq:blend}
\end{equation}
with default $\alpha = 0.5$.
This differs from the pure multiplicative formula in the M.0 framework
($w = w_{\mathrm{base}} \times \mathrm{sim}$) in one important respect:
the $\alpha$ floor prevents the weight from collapsing to zero even for
completely dissimilar tracts.
At $\alpha = 0.5$ and $\mathrm{sim} = 0$, the weight is halved rather than
zeroed.
This is a deliberate choice to prevent METIS from treating economic zone
boundaries as free cuts; geographic boundary length always contributes at
least $\alpha$ of the weight.

The pure M.0 multiplicative formula can be recovered by setting
\texttt{--econ-alpha 0.0} on the CLI.
The practitioner's default of $\alpha = 0.5$ provides a balanced blend that
preserves geographic connectivity while adding economic community signal.

\subsection{Zero-Vector Handling}

The cosine similarity is undefined when either vector is the zero vector
(pure residential tracts with no WAC records).
The implementation handles the three cases:
\begin{itemize}
  \item Both zero: $\mathrm{sim} = 1.0$ (both residential, keep together).
  \item One zero: $\mathrm{sim} = 0.5$ (one residential, one employment center,
    neutral---geographic weight governs).
  \item Neither zero: standard cosine formula.
\end{itemize}
This handling is conservative and community-correct: residential tracts are
their own economic community (bedroom community), not a non-community that can
be arbitrarily cut.

\subsection{M-Track Primary Metric: Within-District Economic Variance}

The proportionality gap (D-seat share minus D-vote share) is reported for
comparison with other B-track papers.
But the M-track primary metric for M.1 is the \emph{mean within-district
standard deviation of \texttt{jobs\_per\_resident}} across all districts in the plan.
Formally, for a $k$-district plan $\{D_1, \ldots, D_k\}$:
\begin{equation}
  \overline{\sigma}_{\mathrm{JPR}}
  = \frac{1}{k} \sum_{i=1}^{k}
    \mathrm{std}\bigl\{\mathrm{JPR}(t) : t \in D_i\bigr\}
  \label{eq:within-var}
\end{equation}
A lower $\overline{\sigma}_{\mathrm{JPR}}$ means that each district contains
tracts with more similar JPR values---economically more coherent districts.
The economic character weights are expected to reduce $\overline{\sigma}_{\mathrm{JPR}}$
by encouraging tracts with similar JPR values (all residential, or all
employment-center) to cluster together rather than being mixed.
Full computation of $\overline{\sigma}_{\mathrm{JPR}}$ for the three test
states is designated Phase~2 (requires post-run tract-level analysis);
the proportionality gap results in Section~\ref{sec:results} are the Phase~1
findings.
